\chapter*{Introduzione}
\markboth{Introduzione}{Introduzione}

L'intelligenza artificiale trova spazio in attività quotidiane e professionali, offrendo strumenti con cui elaborare informazioni e affrontare compiti differenti. Tra questi, i modelli linguistici di grandi dimensioni, indicati con l'acronimo LLM (\textit{Large Language Models}), permettono di interagire attraverso il linguaggio naturale e possono affiancare l'utente nella scrittura, nella sintesi di testi e nell'organizzazione delle idee \cite{Korinek_2023}. Il loro interesse risiede anche nella possibilità di rendere accessibili queste funzioni mediante un'interazione diretta, nella quale una richiesta può essere precisata e sviluppata nel corso del dialogo.

Il valore di tali strumenti dipende dal modo in cui vengono utilizzati e dalla loro adeguatezza al compito. Le risposte devono essere esaminate alla luce delle informazioni disponibili e delle esigenze dell'attività svolta. L'impiego dell'intelligenza artificiale come supporto richiede quindi di comprenderne le capacità e i limiti, mantenendo un ruolo attivo nell'interpretazione dei risultati.

Le applicazioni dell'intelligenza artificiale comprendono anche l'analisi dei dati visivi. La visione artificiale sviluppa metodi per estrarre informazioni dalle immagini e riconoscere caratteristiche, oggetti o variazioni del loro aspetto. Nell'ambito industriale, queste capacità possono supportare il controllo di qualità attraverso l'identificazione di elementi che si discostano dalle caratteristiche attese di un prodotto \cite{Yang_et_al._2021,Cui_et_al._2023}.

La varietà dei materiali rende questo compito articolato. Una superficie può presentare una texture regolare oppure variazioni naturali di colore e tessitura; un oggetto può richiedere di considerare anche la forma e la disposizione delle sue parti. Le caratteristiche che definiscono la conformità dipendono pertanto dal materiale e dal prodotto osservato. Anche le anomalie assumono forme differenti: possono riguardare alterazioni cromatiche, tagli, pieghe, perforazioni o altri cambiamenti locali, con dimensioni e posizioni variabili. Distinguere un difetto da una variazione ammessa richiede una descrizione della normalità coerente con il contesto di osservazione \cite{Cui_et_al._2023}.

Un'ulteriore difficoltà riguarda la disponibilità di esempi. Quando le anomalie sono poco frequenti, raccogliere immagini che rappresentino tutte le possibili tipologie di difetto può risultare difficile. Inoltre, i campioni disponibili non esauriscono necessariamente le manifestazioni che potranno presentarsi in seguito. Questa condizione motiva l'interesse per metodi che apprendono le caratteristiche dei campioni normali e individuano, nelle nuove immagini, deviazioni rispetto alle regolarità apprese \cite{Pang_et_al._2021,Yang_et_al._2021}.

Da queste considerazioni nasce il presente project work, dedicato al rilevamento delle anomalie visive nell'ispezione industriale. Il lavoro confronta quattro modelli, CCW-SSIM, PaDiM, RD4AD e FastFlow, che adottano approcci differenti alla descrizione della normalità. La sperimentazione riguarda la categoria \texttt{leather} del dataset MVTec AD, scelta come caso di studio comune. La varietà dei materiali costituisce il contesto del problema; il confronto sperimentale si concentra invece su un unico materiale, per analizzare il comportamento delle pipeline nelle condizioni adottate.

La domanda che guida l'elaborato riguarda la capacità di questi approcci di riconoscere immagini difettose e di localizzare le regioni anomale. L'esame delle metriche è affiancato dalla lettura delle mappe prodotte, così da collegare i risultati numerici alle risposte osservabili nelle immagini. Il capitolo seguente introduce i concetti necessari a inquadrare il problema e i criteri con cui verrà valutato il confronto.
